\subsection{Commands}

The  head  position  of   the   client's   own   snake  exists  in  two  states
simultaneously: An \emph{acknowledged head}  and  a \emph{predicted head}.

Every  simulation update, a \emph{command} is generated from the current  mouse
and  keyboard  state.  This  command  contains information such as the  snake's
direction and speed.  The  snake's  \emph{predicted head} is stepped using this
command, and the command is then added to the client's ring-buffer. The snake's
\emph{acknowledged head} remains unchanged.

Every  network  update, the entire ring-buffer of commands is  compresssed  and
sent to the server via \verb$MSG_COMMANDS$. This is so that if any commands are
lost, they can be filled in by the next packet again.

The server will  run  its own simulation using the same commands, and send back
an authoritative head position with \verb$MSG_SNAKE_HEAD$.  The client compares
the server's head with its own \emph{acknowledged head}. The  majority  of  the
time, the heads will agree. If  they  differ, then the client rolls its snake's
state back to the frame of the authoritative head, corrects its head so that it
is  identical to the authoritative head, and then replays all commands  in  the
command buffer to step the snake back to the predicated frame.

Visually, when this happens, you will see a slight jitter.

\begin{tikzpicture}[>=latex]
\coordinate (A) at (2,5);
\coordinate (B) at (2,0);
\coordinate (C) at (6,5);
\coordinate (D) at (6,0);
\draw[thick] (A)--(B) (C)--(D);
\draw (A) node[above]{\Large Client};
\draw (C) node[above]{\Large Server};

\coordinate (E) at ($(A)!.1!(B)$);
\coordinate (F) at ($(C)!.25!(D)$);
\coordinate (G) at ($(A)!.39!(B)$);
\coordinate (H) at ($(C)!.35!(D)$);
\coordinate (I) at ($(A)!.6!(B)$);
\draw[->] (E) -- (F) node[midway,sloped,above]{\verb$COMMANDS$};
\draw[->] (H) -- (I) node[midway,sloped,above]{\verb$SNAKE_HEAD$};
\draw (G) node[left]{
    \begin{bytefield}[bitwidth=4em]{1}
        \bitbox{1}{$\text{cmd}_n$} \\
        \bitbox{1}{\ldots} \\
        \bitbox{1}{$\text{cmd}_2$} \\
        \bitbox{1}{$\text{cmd}_1$} \\
    \end{bytefield}};
\draw[|<->|] ([xshift=-50]E) -- ([xshift=-50]I) node[midway,left]{Command buffer};
\end{tikzpicture}



\vspace{1.5em}

\begin{figure}[H]
\verb$MSG_COMMANDS$
\vspace{.25em}\\

\begin{bytefield}[bitwidth=3em,bitformatting=\small]{1}
    \bitheader{0-6} \\
    \bitbox{2}{frame}
    \bitbox{1}{num}
    \bitbox{1}{angle}
    \bitbox{1}{speed}
    \bitbox{1}{action}
    \bitbox{4}{...}
\end{bytefield}

\begin{itemize}
    \item\textbf{frame:} Frame number of the first command in the message.
    \item\textbf{num:} Number of additional commands in the buffer. Can be 0.
    \item\textbf{angle:} 8-bit angle of the snake's  head. Needs to be converted to
    Q4.12 before it can be processed by simulation.
    \item\textbf{speed:} Speed of snake.
    \item\textbf{action:} Action bit-field. For example, if the snake is currently boosting or not.
    \item\textbf{...} \textbf{num} additional commands are delta-compressed and stored here.
\end{itemize}


\end{figure}

\vspace{1.5em}

\begin{figure}[H]
\verb$MSG_SNAKE_HEAD$
\vspace{.25em}\\

\begin{bytefield}[bitwidth=3em,bitformatting=\small]{1}
    \bitheader{0-13} \\
    \bitbox{2}{frame}
    \bitbox{3}{pos x}
    \bitbox{3}{pos y}
    \bitbox{2}{angle}
    \bitbox{1}{speed}
    \bitbox{3}{food}
\end{bytefield}

\begin{itemize}
    \item\textbf{frame:} Frame number of when the head was stepped.
    \item\textbf{pos x:} Q10.14 x coordinate of the snake's head position in world space.
    \item\textbf{pos y:} Q10.14 y coordinate of the snake's head position in world space.
    \item\textbf{angle:} Q4.12 angle of the snake's head.
    \item\textbf{speed:} Current speed of the snake's head.
    \item\textbf{food:} Amount  of  food  eaten  by  the  snake. Internally this is
    stored  as  a  \verb$uint32_t$,  but it's not likely  this  number  will  reach
    anywhere near 16 million.
\end{itemize}



\end{figure}
